%-----------------------------------------------------------------------------------------
% Autor dieser Vorlage:
% Stefan Macke (http://fachinformatiker-anwendungsentwicklung.net)
% Permalink zur Vorlage: http://fiae.link/LaTeXVorlageFIAE
%
% Sämtliche verwendeten Abbildungen, Tabellen und Listings stammen von Dirk Grashorn.
%
% Lizenz: Creative Commons 4.0 Namensnennung - Weitergabe unter gleichen Bedingungen
% -----------------------------------------------------------------------------------------

\documentclass[
	ngerman,
	toc=listof, % Abbildungsverzeichnis sowie Tabellenverzeichnis in das Inhaltsverzeichnis aufnehmen
	toc=bibliography, % Literaturverzeichnis in das Inhaltsverzeichnis aufnehmen
	footnotes=multiple, % Trennen von direkt aufeinander folgenden Fußnoten
	parskip=half, % vertikalen Abstand zwischen Absätzen verwenden anstatt horizontale Einrückung von Folgeabsätzen
	numbers=noendperiod % Den letzten Punkt nach einer Nummerierung entfernen (nach DIN 5008)
]{scrartcl}
\pdfminorversion=5 % erlaubt das Einfügen von pdf-Dateien bis Version 1.7, ohne eine Fehlermeldung zu werfen (keine Garantie für fehlerfreies Einbetten!)
\usepackage[utf8]{inputenc} % muss als erstes eingebunden werden, da Meta/Packages ggfs. Sonderzeichen enthalten

% !TEX root = Projektdokumentation.tex

% Hinweis: der Titel muss zum Inhalt des Projekts passen und den zentralen Inhalt des Projekts deutlich herausstellen
\newcommand{\titel}{Order Comment Tool}
\newcommand{\untertitel}{Tool zu Unterstützung der Produktionsplaner}
\newcommand{\kompletterTitel}{\titel{} -- \untertitel}

\newcommand{\autorName}{Thomas Pöhlmann}
\newcommand{\autorAnschrift}{Schwingstraße 10}
\newcommand{\autorOrt}{68199 Mannheim}

\newcommand{\betriebLogo}{LogoBetrieb.pdf}
\newcommand{\betriebName}{\textsc{Camelot ITLab GmbH}}
\newcommand{\betriebAnschrift}{Theodor-Heuss-Anlage. 12}
\newcommand{\betriebOrt}{68165 Mannheim}

\newcommand{\ausbildungsberuf}{Fachinformatiker für Anwendungsentwicklung}
\newcommand{\betreff}{Dokumentation zur betrieblichen Projektarbeit}
\newcommand{\pruefungstermin}{Winter 2018}
\newcommand{\abgabeOrt}{Mannheim}
\newcommand{\abgabeTermin}{15.12.2018}
 % Metadaten zu diesem Dokument (Autor usw.)
% !TEX root = ../Projektdokumentation.tex

% Anpassung an Landessprache ---------------------------------------------------
\usepackage{babel}

% Umlaute ----------------------------------------------------------------------
%   Umlaute/Sonderzeichen wie äüöß direkt im Quelltext verwenden (CodePage).
%   Erlaubt automatische Trennung von Worten mit Umlauten.
% ------------------------------------------------------------------------------
\usepackage[T1]{fontenc}
\usepackage{textcomp} % Euro-Zeichen etc.

% Schrift ----------------------------------------------------------------------
\usepackage{lmodern} % bessere Fonts
\usepackage{relsize} % Schriftgröße relativ festlegen

% Tabellen ---------------------------------------------------------------------
\PassOptionsToPackage{table}{xcolor}
\usepackage{tabularx}
% für lange Tabellen
\usepackage{longtable}
\usepackage{color, array}
\usepackage{ragged2e}
\usepackage{lscape}
\usepackage{tabu}
%\newcolumntype{w}[1]{>{\raggedleft\hspace{0pt}}p{#1}} % Spaltendefinition rechtsbündig mit definierter Breite

% Grafiken ---------------------------------------------------------------------
\usepackage[dvips,final]{graphicx} % Einbinden von JPG-Grafiken ermöglichen
\usepackage{graphics} % keepaspectratio
\usepackage{floatflt} % zum Umfließen von Bildern
\graphicspath{{Bilder/}} % hier liegen die Bilder des Dokuments

% Sonstiges --------------------------------------------------------------------
\usepackage[titles]{tocloft} % Inhaltsverzeichnis DIN 5008 gerecht einrücken
\usepackage{amsmath,amsfonts} % Befehle aus AMSTeX für mathematische Symbole
\usepackage{enumitem} % anpassbare Enumerates/Itemizes
\usepackage{xspace} % sorgt dafür, dass Leerzeichen hinter parameterlosen Makros nicht als Makroendezeichen interpretiert werden

\usepackage{makeidx} % für Index-Ausgabe mit \printindex
\usepackage[printonlyused]{acronym} % es werden nur benutzte Definitionen aufgelistet

% Einfache Definition der Zeilenabstände und Seitenränder etc.
\usepackage{setspace}
\usepackage{geometry}

% Symbolverzeichnis
\usepackage[intoc]{nomencl}
\let\abbrev\nomenclature
\renewcommand{\nomname}{Abkürzungsverzeichnis}
\setlength{\nomlabelwidth}{.25\hsize}
\renewcommand{\nomlabel}[1]{#1 \dotfill}
\setlength{\nomitemsep}{-\parsep}

\usepackage{varioref} % Elegantere Verweise. „auf der nächsten Seite“
\usepackage{url} % URL verlinken, lange URLs umbrechen etc.

\usepackage{chngcntr} % fortlaufendes Durchnummerieren der Fußnoten
% \usepackage[perpage]{footmisc} % Alternative: Nummerierung der Fußnoten auf jeder Seite neu

\usepackage{ifthen} % bei der Definition eigener Befehle benötigt
\usepackage{todonotes} % definiert u.a. die Befehle \todo und \listoftodos
\usepackage[square]{natbib} % wichtig für korrekte Zitierweise

\usepackage[section]{placeins}

% PDF-Optionen -----------------------------------------------------------------
\usepackage{pdfpages}
\pdfminorversion=5 % erlaubt das Einfügen von pdf-Dateien bis Version 1.7, ohne eine Fehlermeldung zu werfen (keine Garantie für fehlerfreies Einbetten!)
\usepackage[
    bookmarks,
    bookmarksnumbered,
    bookmarksopen=true,
    bookmarksopenlevel=1,
    colorlinks=true,
% diese Farbdefinitionen zeichnen Links im PDF farblich aus
    linkcolor=CamelotBlau, % einfache interne Verknüpfungen
    anchorcolor=CamelotBlau,% Ankertext
    citecolor=CamelotBlau, % Verweise auf Literaturverzeichniseinträge im Text
    filecolor=CamelotBlau, % Verknüpfungen, die lokale Dateien öffnen
    menucolor=CamelotBlau, % Acrobat-Menüpunkte
    urlcolor=CamelotBlau,
% diese Farbdefinitionen sollten für den Druck verwendet werden (alles schwarz)
    %linkcolor=black, % einfache interne Verknüpfungen
    %anchorcolor=black, % Ankertext
    %citecolor=black, % Verweise auf Literaturverzeichniseinträge im Text
    %filecolor=black, % Verknüpfungen, die lokale Dateien öffnen
    %menucolor=black, % Acrobat-Menüpunkte
    %urlcolor=black,
%
    %backref, % Quellen werden zurück auf ihre Zitate verlinkt
    pdftex,
    plainpages=false, % zur korrekten Erstellung der Bookmarks
    pdfpagelabels=true, % zur korrekten Erstellung der Bookmarks
    hypertexnames=false, % zur korrekten Erstellung der Bookmarks
    linktocpage % Seitenzahlen anstatt Text im Inhaltsverzeichnis verlinken
]{hyperref}
% Befehle, die Umlaute ausgeben, führen zu Fehlern, wenn sie hyperref als Optionen übergeben werden
\hypersetup{
    pdftitle={\titel -- \untertitel},
    pdfauthor={\autorName},
    pdfcreator={\autorName},
    pdfsubject={\titel -- \untertitel},
    pdfkeywords={\titel -- \untertitel},
}
 % verwendete Packages
% !TEX root = ../Projektdokumentation.tex

% Seitenränder -----------------------------------------------------------------
\setlength{\topskip}{\ht\strutbox} % behebt Warnung von geometry
\geometry{a4paper,left=20mm,right=20mm,top=25mm,bottom=35mm}

\usepackage[
	automark, % Kapitelangaben in Kopfzeile automatisch erstellen
	headsepline, % Trennlinie unter Kopfzeile
	ilines % Trennlinie linksbündig ausrichten
]{scrpage2}

% Kopf- und Fußzeilen ----------------------------------------------------------
\pagestyle{scrheadings}
% chapterpagestyle gibt es nicht in scrartcl
%\renewcommand{\chapterpagestyle}{scrheadings}
\clearscrheadfoot

% Kopfzeile
\renewcommand{\headfont}{\normalfont} % Schriftform der Kopfzeile
\ihead{\large{\textsc{\titel}}\\ \small{\untertitel} \\[2ex] \textit{\headmark}}
\chead{}
\ohead{\includegraphics[scale=0.2]{\betriebLogo}}
\setlength{\headheight}{15mm} % Höhe der Kopfzeile
%\setheadwidth[0pt]{textwithmarginpar} % Kopfzeile über den Text hinaus verbreitern (falls Logo den Text überdeckt)

% Fußzeile
\ifoot{\autorName}
\cfoot{}
\ofoot{\pagemark}

% Überschriften nach DIN 5008 in einer Fluchtlinie
% ------------------------------------------------------------------------------

% Abstand zwischen Nummerierung und Überschrift definieren
% > Schön wäre hier die dynamische Berechnung des Abstandes in Abhängigkeit
% > der Verschachtelungstiefe des Inhaltsverzeichnisses
\newcommand{\headingSpace}{1.5cm}

% Abschnittsüberschriften im selben Stil wie beim Inhaltsverzeichnis einrücken
\renewcommand*{\othersectionlevelsformat}[3]{
  \makebox[\headingSpace][l]{#3\autodot}
}

% Für die Einrückung wird das Paket tocloft benötigt
%\cftsetindents{chapter}{0.0cm}{\headingSpace}
\cftsetindents{section}{0.0cm}{\headingSpace}
\cftsetindents{subsection}{0.0cm}{\headingSpace}
\cftsetindents{subsubsection}{0.0cm}{\headingSpace}
\cftsetindents{figure}{0.0cm}{\headingSpace}
\cftsetindents{table}{0.0cm}{\headingSpace}


% Allgemeines
% ------------------------------------------------------------------------------

\onehalfspacing % Zeilenabstand 1,5 Zeilen
\frenchspacing % erzeugt ein wenig mehr Platz hinter einem Punkt

% Schusterjungen und Hurenkinder vermeiden
\clubpenalty = 10000
\widowpenalty = 10000
\displaywidowpenalty = 10000

% Quellcode-Ausgabe formatieren
\lstset{numbers=left, numberstyle=\tiny, numbersep=5pt, breaklines=true}
\lstset{emph={square}, emphstyle=\color{red}, emph={[2]root,base}, emphstyle={[2]\color{blue}}}

\counterwithout{footnote}{section} % Fußnoten fortlaufend durchnummerieren
\setcounter{tocdepth}{\subsubsectionlevel} % im Inhaltsverzeichnis werden die Kapitel bis zum Level der subsubsection übernommen
\setcounter{secnumdepth}{\subsubsectionlevel} % Kapitel bis zum Level der subsubsection werden nummeriert

% Aufzählungen anpassen
\renewcommand{\labelenumi}{\arabic{enumi}.}
\renewcommand{\labelenumii}{\arabic{enumi}.\arabic{enumii}.}
\renewcommand{\labelenumiii}{\arabic{enumi}.\arabic{enumii}.\arabic{enumiii}}

% Tabellenfärbung:
\definecolor{heading}{rgb}{0.12,0.19,0.36}
\definecolor{headingfont}{rgb}{1,1,1}
\definecolor{odd}{rgb}{0.9,0.9,0.9}
 % Definitionen zum Aussehen der Seiten
% !TEX root = ../Projektdokumentation.tex

% Abkürzungen, ggfs. mit korrektem Leerraum
\newcommand{\bs}{$\backslash$\xspace}
\newcommand{\bspw}{bspw.\xspace}
\newcommand{\bzw}{bzw.\xspace}
\newcommand{\ca}{ca.\xspace}
\newcommand{\dahe}{\mbox{d.\,h.}\xspace}
\newcommand{\etc}{etc.\xspace}
\newcommand{\eur}[1]{\mbox{#1\,\texteuro}\xspace}
\newcommand{\evtl}{evtl.\xspace}
\newcommand{\ggfs}{ggfs.\xspace}
\newcommand{\Ggfs}{Ggfs.\xspace}
\newcommand{\gqq}[1]{\glqq{}#1\grqq{}}
\newcommand{\inkl}{inkl.\xspace}
\newcommand{\insb}{insb.\xspace}
\newcommand{\ua}{\mbox{u.\,a.}\xspace}
\newcommand{\usw}{usw.\xspace}
\newcommand{\Vgl}{Vgl.\xspace}
\newcommand{\zB}{\mbox{z.\,B.}\xspace}

% Befehle für häufig anfallende Aufgaben
\newcommand{\Abbildung}[1]{\autoref{fig:#1}}
\newcommand{\Anhang}[1]{\appendixname{}~\ref{#1}: \nameref{#1} \vpageref{#1}}
\newcommand{\includegraphicsKeepAspectRatio}[2]{\includegraphics[width=#2\textwidth,height=#2\textheight,keepaspectratio]{#1}}
\newcommand{\Zitat}[2][\empty]{\ifthenelse{\equal{#1}{\empty}}{\citep{#2}}{\citep[#1]{#2}}}
\newcommand{\Autor}[1]{\textsc{#1}} % zum Ausgeben von Autoren
\newcommand{\itemd}[2]{\item{\textbf{#1}}\\{#2}} % erzeugt ein Listenelement mit fetter Überschrift

% fügt Tabellen aus einer TEX-Datei ein
\newcommand{\tabelle}[3] % Parameter: caption, label, file
{\begin{table}[htbp]
\centering
\singlespacing
\input{Tabellen/#3}
\caption{#1}
\label{#2}
\end{table}}

\newcommand{\tabelleAnhang}[1] % Parameter: file
{\begin{center}
\singlespacing
\input{Tabellen/#1}
\end{center}}

\newcommand*{\rowstyle}[1]{% sets the style of the next row
	\gdef\@rowstyle{#1}%
	\@rowstyle\ignorespaces%
}
\newcolumntype{=}{% resets the row style
	>{\gdef\@rowstyle{}}%
}
\newcolumntype{+}{% adds the current row style to the next column
	>{\@rowstyle}%
}

% einfaches Wechseln der Schrift, z.B.: \changefont{cmss}{sbc}{n}
\newcommand{\changefont}[3]{\fontfamily{#1} \fontseries{#2} \fontshape{#3} \selectfont}

% Verwendung analog zu \includegraphics
\newlength{\myx} % Variable zum Speichern der Bildbreite
\newlength{\myy} % Variable zum Speichern der Bildhöhe
\newcommand\includegraphicstotab[2][\relax]{%
% Abspeichern der Bildabmessungen
\settowidth{\myx}{\includegraphics[{#1}]{#2}}%
\settoheight{\myy}{\includegraphics[{#1}]{#2}}%
% das eigentliche Einfügen
\parbox[c][1.1\myy][c]{\myx}{%
\includegraphics[{#1}]{#2}}%
}

\definecolor{CamelotBlau}{rgb}{0.12,0.19,0.36}

% verschiedene Befehle um Wörter semantisch auszuzeichnen ----------------------
\newcommand{\Index}[2][\empty]{\ifthenelse{\equal{#1}{\empty}}{\index{#2}#2}{\index{#1}#2}}
\newcommand{\Fachbegriff}[2][\empty]{\ifthenelse{\equal{#1}{\empty}}{\textit{\Index{#2}}}{\textit{\Index[#1]{#2}}}}
\newcommand{\NeuerBegriff}[2][\empty]{\ifthenelse{\equal{#1}{\empty}}{\textbf{\Index{#2}}}{\textbf{\Index[#1]{#2}}}}

\newcommand{\Ausgabe}[1]{\texttt{#1}}
\newcommand{\Eingabe}[1]{\texttt{#1}}
\newcommand{\Code}[1]{\texttt{#1}}
\newcommand{\Datei}[1]{\texttt{#1}}

\newcommand{\Assembly}[1]{\textsf{#1}}
\newcommand{\Klasse}[1]{\textsf{#1}}
\newcommand{\Methode}[1]{\textsf{#1}}
\newcommand{\Attribut}[1]{\textsf{#1}}

\newcommand{\Datentyp}[1]{\textsf{#1}}
\newcommand{\XMLElement}[1]{\textsf{#1}}
\newcommand{\Webservice}[1]{\textsf{#1}}

\newcommand{\Refactoring}[1]{\Fachbegriff{#1}}
\newcommand{\CodeSmell}[1]{\Fachbegriff{#1}}
\newcommand{\Metrik}[1]{\Fachbegriff{#1}}
\newcommand{\DesignPattern}[1]{\Fachbegriff{#1}}

% Tabellen und Abbildungsverzeichnis
% ------------------------------------------------------------------------------
\renewcommand{\cftfigpresnum}{Abbildung\ }
\renewcommand{\cfttabpresnum}{Tabelle\ }

\newlength{\mylenf}
\settowidth{\mylenf}{\cftfigpresnum}
\setlength{\cftfignumwidth}{\dimexpr\mylenf+2.5em}
\setlength{\cfttabnumwidth}{\dimexpr\mylenf+1em}
 % eigene allgemeine Befehle, die z.B. die Arbeit mit LaTeX erleichtern
% Tabellen, die den Zwischenstand nach einer Projektphase verdeutlichen
\newcommand{\Zwischenstand}[2]{\subsection{Zwischenstand}Tabelle~\ref{tab:#2} zeigt den Zwischenstand nach der #1.\tabelle{Zwischenstand nach der #1}{tab:#2}{#2.tex}}

% Abkürzungen
\newcommand{\Versis}{\textsc{Versis}\xspace}
\newcommand{\NI}{NatInfo\xspace}
\newcommand{\AO}{\textsc{Alte Oldenburger} Krankenversicherung\xspace}
 % eigene projektspezifische Befehle, z.B. Abkürzungen usw.

\begin{document}

\phantomsection
\thispagestyle{empty}
\pdfbookmark[1]{Eidesstattliche Erklärung}{ihkdeckblatt}
\includegraphicsKeepAspectRatio{DeckblattIHK}{1}
\cleardoublepage

\phantomsection
\thispagestyle{plain}
\pdfbookmark[1]{Deckblatt}{deckblatt}
% !TEX root = Projektdokumentation.tex
\begin{titlepage}

\begin{center}
\includegraphics[scale=0.25]{LogoIHK.pdf}\\[1ex]
\Large{Abschlussprüfung \pruefungstermin}\\[3ex]

\Large{\ausbildungsberuf}\\
\LARGE{\betreff}\\[4ex]

\huge{\textbf{\titel}}\\[1.5ex]
\Large{\textbf{\untertitel}}\\[4ex]

\normalsize
Abgabetermin: \abgabeOrt, den \abgabeTermin\\[3em]
\textbf{Prüfungsbewerber:}\\
\autorName\\
\autorAnschrift\\
\autorOrt\\[5ex]

\includegraphics[scale=0.4]{\betriebLogo}\\[2ex]
\textbf{Ausbildungsbetrieb:}\\
\betriebName\\
\betriebAnschrift\\
\betriebOrt\\[5em]
\end{center}

\small
\noindent
Dieses Werk einschließlich seiner Teile ist \textbf{urheberrechtlich geschützt}.
Jede Verwertung außerhalb der engen Grenzen des Urheberrechtgesetzes ist ohne
Zustimmung des Autors unzulässig und strafbar. Das gilt insbesondere für
Vervielfältigungen, Übersetzungen, Mikroverfilmungen sowie die Einspeicherung
und Verarbeitung in elektronischen Systemen.

\end{titlepage}
\cleardoublepage

% Preface --------------------------------------------------------------------
\phantomsection
\pagenumbering{Roman}
\pdfbookmark[1]{Inhaltsverzeichnis}{inhalt}
\tableofcontents
\cleardoublepage

\phantomsection
\listoffigures
\cleardoublepage

\phantomsection
\listoftables
\cleardoublepage

\phantomsection
\lstlistoflistings
\cleardoublepage

\newcommand{\abkvz}{Abkürzungsverzeichnis}
\renewcommand{\nomname}{\abkvz}
\section*{\abkvz}
\markboth{\abkvz}{\abkvz}
\addcontentsline{toc}{section}{\abkvz}
% !TEX root = Projektdokumentation.tex

% Es werden nur die Abkürzungen aufgelistet, die mit \ac definiert und auch benutzt wurden. 
%
% \acro{VERSIS}{Versicherungsinformationssystem\acroextra{ (Bestandsführungssystem)}}
% Ergibt in der Liste: VERSIS Versicherungsinformationssystem (Bestandsführungssystem)
% Im Text aber: \ac{VERSIS} -> Versicherungsinformationssystem (VERSIS)

% Hinweis: allgemein bekannte Abkürzungen wie z.B. bzw. u.a. müssen nicht ins Abkürzungsverzeichnis aufgenommen werden
% Hinweis: allgemein bekannte IT-Begriffe wie Datenbank oder Programmiersprache müssen nicht erläutert werden,
%          aber ggfs. Fachbegriffe aus der Domäne des Prüflings (z.B. Versicherung)

% Die Option (in den eckigen Klammern) enthält das längste Label oder
% einen Platzhalter der die Breite der linken Spalte bestimmt.
\begin{acronym}[WWWWW]
	\acro{ABAP}{Advanced Business Application Programming}
	\acro{ALV}{ABAP List View}
	\acro{APO}{Advanced Planning and Optimization}
	\acro{BAdI}{Business Add-In}
	\acro{ERM}{En\-ti\-ty-Re\-la\-tion\-ship-Mo\-dell}
	\acro{ERP}{Enterprise-Resource-Planning}
	\acro{ECC}{\ac{ERP} Central Component}
	\acro{GUI}{Graphical User Interface}
	\acro{MVC}[MVC]{Model View Controller}
	\acro{PBO}{Process Before Output}
	\acro{PAI}{Process After Input}
	\acro{RFC}{Remote Function Call}
	\acro{UML}{Unified Modeling Language}
\end{acronym}

\clearpage

\section*{Transaktionsverzeichnis}
\markboth{Transaktionsverzeichnis}{Transaktionsverzeichnis}
\addcontentsline{toc}{section}{Transaktionsverzeichnis}
% !TEX root = Projektdokumentation.tex
\textbf{COR1-3}:
\linebreak
Diese Transaktion ist zum Erstellen, bearbeiten und anzeigen von Prozess Aufträgen

\textbf{SAPAPO/RRP3}:
\linebreak
Mithilfe der Produktsicht lassen sich zu einer vorher definierten Produkt-Werks Kombination alle Forcasts, Prozessaufträge, Planaufträge, Produktionsaufträge, Kaufanforderungen und vieles mehr anzeigen lassen.

\textbf{CMOD}:
\linebreak
Mittels dieser Transaktion können Erweiterungspunkte ausgewählt und implementiert werden.

\textbf{se18}:
\linebreak
Diese Transaktion dient zur Implementierung von \ac{BAdI}s.

\textbf{se80}
\linebreak
Die se80 Transaktion ist die am meisten genutzte Transaktion. Mit ihr wird programmiert und es können Datenbanktabellen und Dictionary-Objekte erstellt werden.

\clearpage

% Inhalt ---------------------------------------------------------------------
\pagenumbering{arabic}
% !TEX root = Projektdokumentation.tex
% !TEX root = ../Projektdokumentation.tex
\section{Einleitung}
\label{sec:Einleitung}
Das folgende Projekt ist mein IHK-Abschlussprojekt, welches im Rahmen meiner Ausbildung zum Fachinformatiker im Bereich Anwendungsentwicklung durchgeführt wurde.

\subsection{Vorstellung des Betriebs und meiner Selbst} 
\label{sec:Vorstellung des Betriebs und meiner Selbst}
Ich habe meine Ausbildung am 01.März 2017 bei der Spreitzenbarth Consultants GmbH angefangen. Einem kleinen Unternehmen mit 50 Mitarbeitern mit Fokus auf Supply Chain Management, dort habe ich hauptsächlich Web Anwendungen mit ASP.NET MVC und AngularJs  und Windows Applikationen mit C\# programmiert. Aufgrund einer Firmen Auflösung habe ich zum 01.06.2018 meinen Ausbildungsbetrieb gewechselt und bin zum Camelot ITLab gewechselt. Dort habe ich mich dann mit \ac{ABAP} und dem gesamten SAP Umfeld betätigt.

\subsection{Projektziel} 
\label{sec:Projektziel}
Im SAP \ac{APO} bzw. im SAP \ac{ECC} hat der Planer die Möglichkeit Prozess-, Produktions- und Plan-Aufträge manuell anzulegen bzw. vorhandene Planaufträge zu editieren. Allerdings gibt es hier keine Möglichkeit eine Notiz oder eine Beschreibung zu hinterlegen. Das Projektziel ist die Erstellung von Programmen, die dem Planer ermöglichen, Kommentare	für Aufträge zu erstellen und anzusehen. Außerdem sollen bereits vorhandene Programm entsprechend erweitert werden, um dem Planer den bestmöglichen Komfort zu bieten. Da dieses Projekt ein internes Projekt ist, wurden alle Anforderungen in innerbetrieblichen Besprechungen ausgemacht.  Bei der Implementierung wird besonders darauf geachtet, einen leicht zu erweiternden Code zu produzieren, da später ein Kommentarverlauf hinzugefügt werden soll, sodass der Planer auch vorherige Versionen von Kommentaren anschauen kann. 

\subsection{Projektumfeld}
\label{sec:Projektumfeld}
Die Camelot ITLab GmbH ist eine im SAP-Umfeld tätige Unternehmensberatung, die sowohl funktionale als auch technische Implementierungen von Geschäftsprozessen umsetzt.	Die Abteilung SCM Solution Development innerhalb der Camelot ITLab GmbH, in deren Umfeld	auch dieses Projekt umgesetzt wird, ist für die softwareseitige Implementierung technischer	Anforderungen zuständig. Der Fokus der Abteilung liegt auf Supply Chain Management, im Speziellen Produktions- und Feinplanung.

\subsection{Projektbeteiligte Personen} 
\label{sec:Projektbeteiligte Personen}
Tabelle~\ref{tab:BeteiligtePer} zeigt alle Personen die an dem Projekt beteiligt waren.
\tabelle{Beteiligte Personen}{tab:BeteiligtePer}{BeteiligtePersonen}

\subsection{Projektbegründung} 
\label{sec:Projektbegruendung}
Das Projekt wurde von dem Camelot ITLab angenommen, da es ein häufig von Kunden angefragtes und hilfreiches Feature ist, welches in der standard SAP Lösung fehlt.

//TODO warum Notizen usw.
% !TEX root = ../Projektdokumentation.tex
\section{Projektplanung} 
\label{sec:Projektplanung}

\subsection{Projektphasen}
\label{sec:Projektphasen}
Zur Umsetzung des Projekts standen mir 70 Stunden zur Verfügung. Diese waren sowohl für die Umsetzung des Projekts als auch für die Dokumentation gedacht. Vor Projektbeginn wurde das Projekt in mehrere Phasen gesplittet und die Stunden aufgeteilt. Eine grobe Übersicht kann der nachfolgenden Tabelle~\ref{tab:Zeitplanung} entnommen werden. 

\tabelle{Zeitplanung}{tab:Zeitplanung}{ZeitplanungKurz}
Eine detailliertere Zeitplanung findet sich im \Anhang{app:Zeitplanung}.

\subsection{Ressourcenplanung}
\label{sec:Ressourcenplanung}
In der Ressourcenübersicht, \Anhang{app:Ressourcen}, sind alle Ressourcen aufgelistet, die für das Projekt eingesetzt wurden. Damit sind sowohl Hardware, Software und das Personal gemeint. Es wurde darauf geachtet, dass nur Software zum Einsatz kommt welche entweder kostenfrei (z.B. als Freeware oder gar Open Source) angeboten werden oder die Camelot AG bereits Lizenzen für diese besitzt und zur Verfügung stellen kann. Dadurch sollen die Projektkosten möglichst gering gehalten werden. 

TODO für mehr im anhang schauen ein paar Beispiele.

\subsection{Entwicklungsprozess}
\label{sec:Entwicklungsprozess}
Als Entwicklungsprozess wurde ein agiler Entwicklungsprozess gewählt, sodass während der Implementierung nach jeder Iterationsphase eine Rücksprache mit dem Ausbilder, dem Kunden und der Entwicklungsabteilung bestand. Aufgrund dieser Tatsache wurde bei der Projektplanung relativ wenig Zeit für die Entwurfsphase eingeplant, da sich Teile dieser Phase erst während der Entwicklung ergeben. Die stetige Kommunikation mit der Entwicklungsabteilung fördert außerdem das Erzielen eines besseren Resultats.

% !TEX root = ../Projektdokumentation.tex
\section{Analysephase} 
\label{sec:Analysephase}

\subsection{Ist-Analyse} 
\label{sec:IstAnalyse}
Das Projekt wird im Umfeld der Produktionsplanung durchgeführt und bezieht sich auf alle Zutun Abgangselemente. Der Planer hat im SAP \ac{APO} die	Möglichkeit Zugangs und Abgangselemente manuell anzulegen. In Folgeprozessen kann es zu	Problemen führen, wenn die Gründe dieser Planänderungen nicht sichtbar gemacht werden. Bisher müssen solche Änderungen ohne Systemunterstützung per Email oder auf anderem Wege an alle beteiligten mitgeteilt werden.

\subsection{Wirtschaftlichkeitsanalyse}
\label{sec:Wirtschaftlichkeitsanalyse}
Aufgrund der Probleme, die bereits in der Projektbegründung und der Ist-Analyse beschrieben wurden, ist dieses Projekt absolut notwendig. Trotzdem werde ich in dem folgenden Abschnitt analysieren ob die Umsetzung auch aus Wirtschaftlichen Gesichtspunkten gerechtfertigt ist.

\subsubsection{\gqq{Make or Buy}-Entscheidung}
\label{sec:MakeOrBuyEntscheidung}

Da es etwas vergleichbares auf dem Markt noch nicht gibt und das Camelot ITLab davon lebt Software zu produzieren und zu verkaufen, würde es natürlich nix bringen die Software zu kaufen. Allerdings stellt sich die Frage selbst produzieren oder von einer fremd Firma oder Freelancern im Namen des Camelot ITLabs entwickeln zu lassen. In der Folgende Tabelle werden alle ausschlaggebenden Punkte welche für die jeweilige Option sprechen aufgeführt.

Camelot Style und programming rechtlinien								Bei Problemen oder nicht einhaltung des Zeitplans
																	    ist der andere Schuld
Support nach der fertigstellung muss geklärt werden

Kosten rechnung

Aufgrund der oben genannten Punkte wurde entschieden die Software selber zu produzieren. 
  

Irwelche Inder kosten  * Zeit.
Später support schwierig da man entweder die Entwickler halten muss oder sie halt weggehen lässt.

Unkalkulierbare kosten nach Projektabschluss.
Da wir davon leben Produkte zu entwickeln und zu lizensieren. 

\subsubsection{Projektkosten}
\label{sec:Projektkosten}
Die Projektkosten, die während dem gesamten Entwicklungszeitraum anfangen sollen im Folgenden kalkuliert werden. Dafür müssen neben den Kosten, die für Hardware und Software anfallen auch die Personalkosten berücksichtigt werden. Da die genauen Personalkosten Betriebsgeheimnis sind, wird die Kalkulation mit groben Stundensätzen durchgeführt. Der Stundensatz eines Auszubildenden ergibt sich aus dem Monats ca. 1000€.
		
\begin{eqnarray}
	365 \mbox{ Tage/Jahr} \cdot \frac{5}{7} - 30 Tage = 260 \mbox{ Tage/Jahr}\\
	8 \mbox{ h/Tag} \cdot 230 \mbox{ Tage/Jahr} = 1840 \mbox{ h/Jahr}\\
	\eur{1000}\mbox{/Monat} \cdot 12 \mbox{ Monate/Jahr} = \eur{12000} \mbox{/Jahr}\\
	\frac{\eur{12000} \mbox{/Jahr}}{1840 \mbox{ h/Jahr}} \approx \eur{6,52}\mbox{/h}
\end{eqnarray}
	
Der Rechnung zufolge beträgt der Stundensatz eines Azubis ca. 7€ die Stunde und die kosten eines normalen Mitarbeiters schätze ich auf 25€ die Stunde. Für die Ressourcen werden die gesamten Hardwarekosten addiert und dann durch die durchschnittliche Lebenszeit in Arbeitsstunden, d.h. 3 * 365 * 5/7 * 8 geteilt.
	
Tabelle~\ref{tab:HardwareKosten} zeigt die Hardwarekosten, die für die Geräte eines Mitarbeiters anfallen.
\tabelle{HardwareKosten}{tab:HardwareKosten}{HardwareKosten}\\
	
Daraus Ergibt sich dann folgende Rechnung und Stundesatz
\begin{eqnarray}
	\frac{\eur{1680} \mbox{/3 Jahre}}{3} = 560 \mbox{/Jahr}\\
	\frac{\eur{560} \mbox{/Jahre}}{1840 \mbox{ h/Jahr}} \approx \eur{0.30}\mbox{/h}
\end{eqnarray}

Es Ergibt sich also ein Stundensatz von 0.3€ für die Hardware Kosten. Dazu kommen natürlich noch Kosten für Arbeitsplatz, Möblierung, Strom und Internet.
\tabelle{Projektkosten}{tab:Projektkosten}{Projektkosten.tex}

\subsubsection{Amortisationsdauer}
\label{sec:Amortisationsdauer}
Im folgenden Abschnitt soll berechnet werden, ab welchem Zeitraum sich die Entwicklung auch aus Wirtschaftlichem Sichtpunk für das Camelot ITLab und für den Kunden lohnt. Die Amortisationsdauer wird berechnet indem man die Produktionskosten bzw. Anschaffungskosten durch die Gewinne dividiert, welche durch das neue Produkt erzielt wurden.

\paragraph{Camelot ITLab}
Die Gewinne welche durch dieses Produkt entstehen lassen sich in zwei Kategorien teilen. Zum einen die Gewinne welche durch den tatsächlichen verkauf dieses Produkt erzielt werden und zum anderen die Gewinne die erzielt werden, wenn der Kunde außerdem neben diesem Tool dann auch noch andere Software der Camelot kaufen möchte. Für diese Tool wird ein Preis von etwa 500€ angesetz. Der Break-even-point (Gewinnschwell) liegt bei
\begin{eqnarray}
\frac{\eur{1510}}{\eur{500}} \approx 3 
\end{eqnarray}

\paragraph{Kunde}
In Fall dieses Projektes lassen sich die Gewinne nicht so leicht erfassen, da zum einen die Fehleranfälligkeit der alten Lösung, Kommentare per Email verschicken, Zettelwirtschaft, usw. deutlich verringert wird. Zum Anderen erspart das Tool dem Planer, welcher die Prozess anlegt deutlich Zeit, da dieser nicht noch ein weiteres Programm benötigt um die Kommentare einzutragen und an seine Kollegen zu verteilen und den Mitarbeitern an der Maschine Zeit, welche die Kommentare direkt in der Transaktion neben den Aufträgen sehen deutlich Zeit. 

\paragraph{Beispielrechnung (verkürzt)}
Es wird davon ausgegangen, dass ein Planer jeden Tag ca. eine Stunde damit beschäftigt ist Aufträge zu Kommentieren. Dieses Tool bietet eine Zeitersparung von ca. 50\% da die Texte weiterhin per Hand verfasst werden müssen. Dies bedeutet Zeiteinsparung von 30 Minuten pro Tag und pro Planer. Bei etwa 230 Arbeitstagen pro Jahr ergibt sich eine gesamte Zeiteinsparung von  
\begin{eqnarray}
230 \mbox{ Tage/Jahr} \cdot 30 \mbox{ min/Tag} = 6900 \mbox{ min/Jahr} = 115 \mbox{ h/Jahr} 
\end{eqnarray}

Das Gehalt eines Produktionsplaners beträgt laut Experten etwa 14€ die Stunde. Die Tatsächlichen kosten, die für die Firma anfallen, werden auf ca. 20€ geschätzt.
Dadurch ergibt sich eine jährliche Einsparung von 
\begin{eqnarray}
115 \mbox{h} \cdot \eur{(20)}{\mbox{/h}} = \eur{2300}
\end{eqnarray}

Je nach Firmengröße und Anzahl der Planer variiert die Amortisationszeit stark. Daher wurden Anhand dreier Beispiele die Zeit berechnet.

Die Amortisationszeit für eine kleine Firma mit 2 Planern beträgt:
$\frac{\eur{500}}{2 \cdot \eur{2300€} \mbox{/Jahr} } \approx 0,19 \mbox{ Jahre} \approx  \mbox{ Wochen}$.

Die Amortisationszeit für eine mittelgroße Firma mit 50 Planern beträgt
$\frac{\eur{500}}{50 \cdot \eur{2300€} \mbox{/Jahr} } \approx 0,004 \mbox{ Jahre} \approx 1,5 \mbox{ Tage}$.

Die Amortisationszeit für eine große Firma mit 200 Planern beträgt
$\frac{\eur{500}}{200 \cdot \eur{2300€} \mbox{/Jahr} } \approx 0,001 \mbox{ Jahre} \approx 9 \mbox{ Stunden}$.

\subsection{Nutzwertanalyse}
\label{sec:Nutzwertanalyse}
Lohnt sich das Projekt für den Kunden was wir gespart weniger fehler usw. alles außer geld


\subsection{Anwendungsfälle}
\label{sec:Anwendungsfaelle}

An der Maschiene kann der arbeiter direkt neben dem auftrag auch den Kommentar sehen.

Use-Case Diagram

\subsection{Qualitätsanforderungen}
\label{sec:Qualitaetsanforderungen}
Wie in jedem anderen Projekt des Camelot ITLabs gelten auch hier die programming Guidelines. Ein kleine Zusammenfassung mit allen für dieses Projekt wichtigen Punkten findet sich im Anhang \Anhang{app:Programming Guidlines}. Außerdem muss es für den Administration Screen einen Eingabeüberprüfung geben, der verhindert, dass der User Keyfelder bzw. Fleder die es nicht gibt angibt.

\subsection{Lastenheft/Fachkonzept}
\label{sec:Lastenheft}
\begin{itemize}
	\item Auszüge aus dem Lastenheft/Fachkonzept, wenn es im Rahmen des Projekts erstellt wurde.
	\item Mögliche Inhalte: Funktionen des Programms (Muss/Soll/Wunsch), User Stories, Benutzerrollen
\end{itemize}

\paragraph{Beispiel}
Ein Beispiel für ein Lastenheft findet sich im \Anhang{app:Lastenheft}. 

\Zwischenstand{Analysephase}{Analyse}

% !TEX root = ../Projektdokumentation.tex
\section{Entwurfsphase} 
\label{sec:Entwurfsphase}

\subsection{Zielplattform}
\label{sec:Zielplattform}
Das Abschlussprojekt soll wie bereits im Projektziel beschrieben eine Erweiterung zu bereits vorhandenen Transaktionen darstellen als auch eigene Programme zur Verwaltung und Massenpflege der Kommentare mit sich bringen. Als Programmiersprache wird ABAP (Advanced Bussiness Application Programming) verwendet, eine eigens von der SAP entwickelte Programmiersprache, die in ihrer Grundstruktur der Sprache COBOL ähnelt.

\subsection{Architekturdesign}
\label{sec:Architekturdesign}
Die Programme im \ac*{APO} und \ac*{ECC} werden nach dem \ac*{MVC} Konzept programmiert. Allerdings habe ich in diesem Fall kein Model benötigt. Es gibt in jedem System jeweils eine \ac*{GUI} Klasse, welche für die visuelle Darstellung und die Reaktion auf Benutzerinteraktionen zuständig ist. Diese Klasse besitzt für jeden Screen eine Member Struktur, die die anzuzeigenden Daten hält und jeweils eine \ac*{PBO} und \ac*{PAI} Methode. Die jeweiligen Screens werden in einer Funktionsgruppe definiert und mithilfe des SAP Screen Panters gestaltet. Außerdem gibt es jeweils eine Controller Klasse, welcher für die gesamte Logik zuständig ist. Die Startpunkte der Programme sind jeweils ein Report welcher, entweder einen Selektionsbildschirm hat oder direkt über ein Funktionsmodul, welches in der Funktionsgruppe definiert ist den gewünschten screen startet, da aus einem Report direkt kein Screen aufgerufen werden kann. Das \ac*{PBO} und das \ac*{PAI} in der Funktionsgruppe sind dynamisch agierende Module, welche dann auf die jeweilige Methode der \ac*{GUI} Klasse weiterleiten.

\subsection{Entwurf des Userinterface}
\label{sec:Benutzeroberflaeche} 
Um die Anwendungen möglichst Benutzerfreundlich zu gestalten wurde im Vorfeld klar strukturiert auf welchem Screen der User welche Informationen angezeigt bekommen soll. Außerdem wurden die Möglichen Selektionskriterien vorab geplant damit später keine Zeit mit der Erstellung von unnötigen Datenstrukturen verbraucht wird. Es wird später im \ac*{ECC} einen Screen mit dem Namen Administration geben, auf welchem der Planer ein Feld der Datenbank Tabelle AUFK angeben kann, in welchem dann der Kommentar gespeichert wird. Hier wird darauf geachtet, dass der code leicht zu erweitern ist, da hier später auch noch weitere Tabellen und Felder ausgewählt werden sollen können wie z.B. die PLAF, in welcher Planaufträge gespeichert sind. Auf dem Hauptscreen des Programms werden die Auftrags Daten und den Kommentaren in einem ALV (ABAP List View) dargestellt. Im APO wird es denselben Hauptscreen auch geben hier fällt allerdings der Administration Screen weg, da hier alle Kommentare unabhängig der Kategorie in der Datenbank /SAPAPO/ORDFLDS gespeichert werden. 

\subsection{Datenmodell}
\label{sec:Datenmodell}
Im ECC werden zwei Datenbank Tabellen erstellt. Zum einen die /CAMELOT/OC\_COMT. Die Beziehungen sind in dem Folgenden \ac*{UML}.

\begin{figure}[htb]
\centering
\includegraphicsKeepAspectRatio{DatabaseTablesECC01.pdf}{0.7}
\caption{Vereinfachtes ER-Modell}
\label{fig:ECC01}
\end{figure} 

Die Tabelle AUFK wird mittels einem Custom Include um ein Feld ZZ\_ORDER\_COMMENT erweitert. Dieses wird für das COR1 Enhancement gebraucht. Außerdem wird eine Datenbank Tabelle /CAMELOT/OC\_SETT erstellt, wie in dem folgenden \ac*{UML} beschrieben. In dieser Datenbank wird vorerst lediglich ein Feldname der AUFK gespeichert. Später muss diese Tabelle dann noch um Tabellenname erweitert werden.

\begin{figure}[htb]
	\centering
	\includegraphicsKeepAspectRatio{DatabaseTablesECC02.pdf}{0.2}
	\caption{Vereinfachtes ER-Modell}
	\label{fig:ECC02}
\end{figure} 

Im APO benötigen wir keine Extra Datenbank zum Speichern der Kommentare, da wir hier, wie bereits oben erwähnt, die /SAPAPO/ORDFLDS per Custom Include um ein FELD ORDER\_COMMENT erweitern können.  

\begin{figure}[htb]
	\centering
	\includegraphicsKeepAspectRatio{DatabaseTablesAPO01.pdf}{0.2}
	\caption{Vereinfachtes ER-Modell}
	\label{fig:APO01}
\end{figure} 

\subsection{Maßnahmen zur Qualitätssicherung}
\label{sec:Qualitaetssicherung}
Um die hohen Qualitätsanforderungen der Camelot ITLab zu gewährleisten und die Qualitätsanforderungen des Projekts zu sicher, werden während der laufenden Entwicklung nach jedem Iterationsschritt die neu eingebauten Funktionalitäten getestet. Außerdem wird es mehrere Code Review geben in denen andere Entwickler sich den Code anschauen und gegebenenfalls Schwachstellen erkennen und Verbesserungsvorschläge einbringen. Alle Tests werden manuell durchgeführt.  


\subsection{Pflichtenheft/Datenverarbeitungskonzept}
\label{sec:Pflichtenheft}
\begin{itemize}
	\item Auszüge aus dem Pflichtenheft/Datenverarbeitungskonzept, wenn es im Rahmen des Projekts erstellt wurde.
\end{itemize}

\paragraph{Beispiel}
Ein Beispiel für das auf dem Lastenheft (siehe Kapitel~\ref{sec:Lastenheft}: \nameref{sec:Lastenheft}) aufbauende Pflichtenheft ist im \Anhang{app:Pflichtenheft} zu finden.


\Zwischenstand{Entwurfsphase}{Entwurf}

% !TEX root = ../Projektdokumentation.tex
\section{Implementierungsphase} 
\label{sec:Implementierungsphase}

\subsection{Iterationsplanung}
\label{sec:Iterationsplanung}
Bevor mit der eigentlichen Implementierung begonnen wurde, wurde zuerst ein Interationsplan erstellt. In ihm werden die Iterationsschritte und deren Reihenfolge definiert. innerhalb einer Iteration wird die zuvor definierte Funktionalität eingebaut. Der erstellte Iterationsplan befindet sich im Anhang \Anhang{app:Iterationsplan}.

\subsection{Implementierung der Datenstruktur ECC}
\label{sec:Implementierung der Datenstruktur ECC}
Zuerst wurden alle Dictionary-Objekte, welche im \ac{ECC} gebraucht werden, erstellt. Eine Vollständige Liste aller Tabellen Typen, Strukturen, Datenelemente und Domänen können dem Anhang entnommen werden. Außerdem wurde die Datenbank Tabellen wie in Kapitel~\ref{sec:Datenmodell}: ~\nameref{sec:Datenmodell} beschrieben implementiert. Die AUFK wurde mittels einem Custom Include um das Feld ZZ\_ORDER\_COMMENT erweitert.

\subsection{Implementierung der Benutzeroberfläche ECC}
\label{sec:Implementierung der Benutzeroberfläche ECC}
In der \ac{GUI} Funktionsgruppe wurden jeweils ein Screen für das Maintenance Programm (Screen 0100) und ein Screen für das Administration Programm (Screen 0500) mithilfe des Screen Panters angelegt und gestaltet. Der Screen 0100 enthält lediglich einen Custom Container, in welchem dann später das \ac{ALV} angezeigt wird. Der Screen 0500 hat zum jetzigen Zeitpunkt nur ein Label und ein Eingabefeld, dessen Element sich in der Member Struktur des Screens der \ac{GUI} Klasse befindet. Das \ac{PBO} und das \ac{PAI} der Funktionsgruppe wurde dynamisch programmiert. Das bedeutet, dass je nach Screen die zugehörige Methode in der \ac{GUI} Klasse aufgerufen wird. Die User Befehle (OK Code) werden in ein Member Feld der \ac{GUI} Klasse geschrieben. Außerdem wurde für das Maintenance Programm ein extra Selektionsbildschirm angelegt, in welchem man gewisse Parameter wie Benutzer, Auftragsnummer usw. eingeben kann und somit die dargestellten Aufträge gefiltert werden. Screenshots der Anwendung befinden sich unter \Anhang{Screenshots}.

\paragraph{Beispiel}
Screenshots der Anwendung in der Entwicklungsphase mit Dummy-Daten befinden sich im \Anhang{Screenshots}.

\subsection{Implementierung PBO und PAI}
\label{sec:Implementierung PBO und PAI}
Jeder Screen der Funktionsgruppe hat seine eigene \ac{PBO} und \ac{PAI} Methode in der GUI Klasse. Im PBO werden die Daten geladen bevor sie dann auf dem Screen angezeigt werden. Dies passiert, indem die jeweilige Controller Methode aufgerufen wird. Die Daten werden allerdings nicht jedes Mal neu geladen, wenn wir das \ac{PBO} betreten, sondern nur, wenn die ALV Struktur, welche in der Member Struktur des jeweiligen Screens ist, Initial also "leer" ist. Dadurch wird verhindert das wir zu viele Datenbank Zugriffe haben, da jedes Mal die Daten neu geladen werden. Nichts desto trotz können wir einen Refresh einleiten indem wir einfach die \ac{ALV} Struktur Clearen. Im \ac{PAI} werden all User Befehle (OK Code) abgefangen welche vom Screen geworfen werden und dann die jeweilige Action ausgeführt. 

\subsection{Implementierung der Geschäftslogik ECC}
\label{sec:Implementierung der Geschäftslogik ECC}
Die eigentliche Geschäftslogik und Datenbank zugriffe finden alle im Controller statt. Hier wurden mehrere Methoden implementiert, welche die gesamten Aufträge, die derzeit unterstützt werden, (Prozessaufträge, Produktionsaufträge aus der AUFK mit join auf die AFKO für Material Infos und Planaufträge aus der PLAF) laden. Außerdem wurde die Speicherlogik implementiert. In der Member Struktur des Screens gibt es zwei Tabellen eine für die Daten, die dann tatsächlich im \ac{ALV} angezeigt werden und eine andere in welcher immer die originalen Daten seit dem letzten speichern drin sind. Beim speichern werden nun diese zwei Tabellen verglichen und so alle Aufträge, welche sich nicht geändert haben aussortiert, sodass die Speicherlogik nur auf die tatsächlich modifizierten oder erstellten Aufträge angewendet wird. Diese Aufträge werden dann mithilfe eines RFC ins \ac{APO} System transferiert. (siehe Kapitel~\ref{sec:Implementierung der Geschäftslogik APO}: \nameref{sec:Implementierung der Geschäftslogik APO})     

\subsection{Implementierung der COR Erweiterung}
\label{sec:Implementierung der COR Erweiterung}   
Um die COR Transaktion (1-3) zu erweitern musste zunächst einmal ein passender Erweiterungspunkt gefunden werden, welcher zum einen, einen Screen Exit und zum anderen, zwei Functions Exits vor und nach dem laden der Daten hat. Verwendet wurde das Enhancement PPCO0020. Diese bietet alle Komponenten, die ich zum anzeigen der Daten benötige. Zuerst wurde das Screen Exit erstellt und mittels des Screen Painters ein simpler Screen mit einem Label und einem Eingabefeld erstellt und aktiviert. Dann wurde das \ac{PBO} Modul um eine Methode erweitert, sodass das Eingabefeld in der COR3 deaktiviert wird, da diese Transaktion nur zum Anzeigen ist. Außerdem wurde in dem Top Include eine globale Struktur vom typ AUFK angelegt, dessen ZZ\_ORDER\_COMMENT Feld hinter dem Eingabe Feld liegt. In dem ersten Function Exit, wird diese Struktur dann gefüllt. Da der Planer aber auch die Möglichkeit haben soll Kommentar von Prozessaufträgen zu ändern bzw. beim anlegen eines Auftrags anzugeben wird noch ein weiterer Funktionsbaustein benötigt, da der oben genannte, keinen Function Exit für das Speichern hat, von wo wir unsere eigene Speicherlogik aus aufrufen können. Hierzu wurde der Erweiterungspunkt PPCO0007 gewählt. Dieser Erweiterungspunkt liefert einen Function Exit aus welchem dann eine statische Methode aus dem Controller aufgerufen wird damit der geänderte Kommentar nicht nur in der AUFK sondern auch in der /CAMELOT/OC\_COMT gespeichert wird und von wo außerdem der \ac{RFC} ausgeführt wird damit der Kommentar ins \ac{APO} transferiert wird. Näheres (siehe Kapitel~\ref{sec:Implementierung der Geschäftslogik APO}: \nameref{sec:Implementierung der Geschäftslogik APO})

\subsection{Implementierung der Datenstruktur APO}
\label{sec:Implementierung der Datenstruktur APO} 
Nachdem im \ac{ECC} System nun alles Implementiert wurde, muss nun ein großteil der selben Datenstrukturen auch im \ac{APO} angelegt werden. Eine Liste aller Dictionary-Objekte kann dem Anhang entnommen werden. Des Weiteren wurde die /SAPAPO/ORDFLDS Datenbank Tabelle um ein Feld, ORDER\_COMMENT, erweitert. Im \ac{APO} wird keine extra Datenbank wie im \ac{ECC} benötigt, da für alle Aufträge die /SAPAPO/ORDFLDS Tabelle benutzt werden kann.

\paragraph{Beispiel}
Screenshots der Anwendung in der Entwicklungsphase mit Dummy-Daten befinden sich im \Anhang{Screenshots}.

\subsection{Implementierung der Benutzeroberfläche APO}
\label{sec:Implementierung der Benutzeroberfläche APO} 
Ähnlich wie im \ac{ECC} musste ein Screen(0100) erstellt werden, auf welchem das \ac{ALV} mit den Auftragsdaten und dem Kommentar angezeigt wird. Der zweite Screen, welchem im \ac{ECC} angelegt wurde, wird im \ac{APO} nicht benötigt, da automatisch alle Kommentaren in der /SAPAPO/ORDFLDS im Feld ORDER\_COMMENT gespeichert werden. Der Screen 0100 ist genauso aufgebaut wie im \ac{ECC} nur der Selektionsbildschirm fällt zum jetzigen Zeitpunkt deutlich kleiner aus, da man nur nach Order Number filtern kann.  

\subsection{Implementierung der Geschäftslogik APO}
\label{sec:Implementierung der Geschäftslogik APO}  
Wie im \ac{ECC} findet auch im \ac{APO} die gesamte Geschäftlogik im Controller statt. Neben den selben Klassen welche es auch im \ac{ECC} gibt(Constants, Controller, \ac{GUI}, \ac{ALV}) gibt es eine extra Klasse mit dem Namen /CAMELOT/CL\_OC\_RRP, welche für den \ac{BAdI} und die implizite Erweiterung der RRP3 genutzt wird. Außerdem gibt es neben der Funktionsgruppe welche für das Hauptprogramm genutzt wird eine weiter Funktionsgruppe /CAMELOT/OC\_COMMENT mit der Funktionsmodul /CAMELOT/OC\_COMMENT\_SYNC welche Remote-Enabled ist. Das bedeutet sie kann aus einem anderen System heraus aufgerufen werden. Diese Methode hat als Importparamter eine Tabelle mit dem Tabellentyp /CAMELOT/OC-\_ORD\_COMMENT\_RFC\_T (siehe \Anhang{Dictionary-Objekte}). Diese Methode Funktion dient zum transport der Kommentare vom \ac{ECC} ins \ac{APO}. Die Strukture der Tabelle beinhaltet nur die Auftragsnummer und den Kommentar um den Traffic möglichst gering zu halten. Die Methode selber sucht mithilfe eines SAP Funktionsbausteins /SAPAPO/DM\_ORDER-\_GET\_ORDID die richtige Orderid aus dem Livecache zu der gegebenen Auftragsnummer. Mithilfe der Orderid und dem Kommentar wird dann die /SAPAPO/ORDFLDS Tabelle geupdatet.  

\subsection{Implementierung der RRP3 Erweiterung}
\label{sec:Implementierung der RRP3 Erweiterung}   
Um die RRP3 zu erweitern wurde nicht wie bei der COR1-3, Customer-Exits, sondern ein \ac{BAdI} und eine implizite Erweiterung verwendet. Der Unterschied zwischen einem Customer-Exits und einem \ac{BAdI} ist, dass der \ac{BAdI} die neuere Objektorientierte Variante eines Customer-Exits ist. Statt in Funktionsmodule Code einzufügen wird eine vordefinierte Klasse implementiert deren Methoden wie bei den alten Customer-Exits zu einem bestimmten Zeitpunkt aufgerufen werden. Für diese RRP3 Erweiterung wurden der \ac{BAdI} /SAPAPO/RRP\_IO\_COL benötigt.Mit diesem war es möglich an den Feldkatalog der RRP3 das Feld Order\_Comment anzufügen und dieses dann auch mittels einer zweiten Methode mit den richtigen Kommentaren zu füllen. Außerdem musste eine Implicites Enhancement angelegt werden, da es mithilfe von Customer-Exits oder \ac{BAdI}s nicht möglich war die t-style Tabelle der einzelnen Elemente in der \ac{ALV} Tabelle zu ändern, was allerding absolut notwendig ist um z.B. den Kommentar editierbar zu machen. Die jeweiligen Methoden des \ac{BAdI}s und der impliziten Erweiterung rufen eine statische Methode in der RRP3 Klasse auf, welche genau den selben Namen hat wie die eigentliche Methode des \ac{BAdI}s, um einen besseren Überblick zu verschaffen und sie später gut bei Kunden implementieren zu können. Diese Methoden holen sich dann die Instance des Controllers bzw. eine wird erstellt und rufen dann die jeweiligen Controller Methoden auf.

\Zwischenstand{Implementierungsphase}{Implementierung}

\input{Inhalt/Qualitätssicherung}
% !TEX root = ../Projektdokumentation.tex
\section{Fazit} 
\label{sec:Fazit}

\subsection{Soll-/Ist-Vergleich}
\label{sec:SollIstVergleich}

todo es hat geklappt

Die Zeit von 70 Stunden und der in Abschnitt ~\ref{sec:Projektphasen} (\nameref{sec:Projektphasen}) erstellte Projektplan wurden insgesamt eingehalten auch wenn es bei den einzelnen Projektphasen zu Verschiebungen gekommen ist. Die Implementierung im \ac{ECC} System hat etwas länger gedauert als zuerst angenommen, dafür konnte ein Großteil der Logik ins \ac{APO} System kopiert und hier einige Zeit gespart werden. In der Tabelle ~\ref{tab:Soll-/Ist-Vergleich}: ~\nameref{tab:Soll-/Ist-Vergleich} wird die tatsächlich benötigte Zeit gegenüber die geplante Zeit gestellt und verglichen.
\tabelle{Soll-/Ist-Vergleich der Projektphasen Zeitplanung}{tab:Soll-/Ist-Vergleich}{Zeitnachher}

\subsection{Ausblick}
\label{sec:Ausblick}
Alle im Lastenheft definierten Anforderungen konnten erfüllt werden, nichts desto trotz möchte ich hier einen kleinen Ausblick über mögliche weiter Feature geben die während der Entwicklung aufgefallen und als praktisch erachtet wurden. Zum einen gibt es noch weitere Transaktionen die erweiter werden können wie z.B. die Feinplantafel um den Planer noch mehr Komfort zu bieten. Außerdem soll man in einer späteren Version im \ac{ECC} im Administration Screen nicht nur einen Feldnamen angeben können, sondern auch die Datenbanktabelle bzw. mehrere Einträge vornehmen können. Außerdem soll in absehbarer zukunft ein Lösch Programm geschrieben werden, dass die Aufträge welche sich im ECC System in der /CAMELOT/OC\_COMT Datenbanktabelle befinden überprüft und Aufträge welche nicht mehr existieren aus der Tabelle löscht. Aufgrund des in Abschnitt ~\ref{sec:Architecturedesign} beschriebenen Aufbaus des Programms lassen sich Änderungen und Anpassungen sehr einfach vornehmen. Außerdem erhöht dies die Wartbarkeit der Programme, da alle Mitarbeiter des Camelot ITLabs sich besondern mit dieser Struktur auskennen. 


% Literatur ------------------------------------------------------------------
\clearpage
\nocite{*}
\renewcommand{\refname}{Literaturverzeichnis}
\bibliography{Bibliographie}
\bibliographystyle{Allgemein/natdin} % DIN-Stil des Literaturverzeichnisses
% !TEX root = Projektdokumentation.tex
\clearpage
\addsec{Eidesstattliche Erklärung}

% Hinweis: die eidesstattliche Erklärung ist ggfs. an die Richtlinie der IHK anzupassen

Ich, \autorName, versichere hiermit, dass ich meine \textbf{\betreff} mit dem
Thema
\begin{quote}
\textit{\kompletterTitel}
\end{quote}
selbständig verfasst und keine anderen als die angegebenen Quellen und Hilfsmittel benutzt habe,
wobei ich alle wörtlichen und sinngemäßen Zitate als solche gekennzeichnet habe. Die Arbeit
wurde bisher keiner anderen Prüfungsbehörde vorgelegt und auch nicht veröffentlicht.\\[6ex]

\abgabeOrt, den \abgabeTermin


\rule[-0.2cm]{5.5cm}{0.5pt}

\textsc{\autorName}


% Anhang ---------------------------------------------------------------------
\clearpage
\appendix
\pagenumbering{roman}
% !TEX root = Projektdokumentation.tex
\section{Anhang}
\subsection{Detaillierte Zeitplanung}
\label{app:Zeitplanung}

\tabelleAnhang{ZeitplanungKomplett}

\subsection{Lastenheft (Auszug)}
\label{app:Lastenheft}

Die Anwendungen müssen folgende Anforderungen erfüllen: 
\begin{enumerate}[itemsep=0em,partopsep=0em,parsep=0em,topsep=0em]
\item Darstellung der Daten
	\begin{enumerate}
	\item Die Anwendung muss die Funktion bieten, Aufträge nach vorgegebenen Kriterien filtern zu können.
	\item Die Aufträge müssen in einer für den Planer übersichtlichen Umgebung dargestellt werden.
	\item Der Planer soll möglichst einfach und schnell Kommentare schreiben und sich anzeigen lassen können. 
	\item Der Planer soll zudem auch in den gewohnten Standardanwendungen Kommentare verfassen und lesen können 
	\end{enumerate}
\item Weitere Anforderungen
	\begin{enumerate}
	\item Die Anwendungen sollen sauber und übersichtlich programmiert werden, sodass spätere Erweiterungen leicht durchzuführen sind.
	\item Möglichst wenig Code soll in User Exits oder sonstigen Erweiterungspunkten geschrieben werden, da diese immer manuell beim Kunden eingefügt werden müssen.
	\end{enumerate}
\end{enumerate}


\clearpage

\subsection{Iterationsplan}
\label{app:Iterationsplan}

% !TEX root = Projektdokumentation.tex
\begin{enumerate}[itemsep=0em,partopsep=0em,parsep=0em,topsep=0em]
	\item ECE Entwicklung
	\begin{enumerate}
		\item Erstellung der Dictionary-Objekte
		\item Erstellung der Datenbanken
		\item Erstellung der Klassen ALV, GUI, Controller und Constants
		\item GUI Funktionsgruppe erstellen
		\item Maintenance Screen erstellen
		\item Administration Screen erstellen
		\item AUFK erweitern und Speicher Logik einbauen
		\item COR Programm erweitern
	\end{enumerate}
	\item SCE Entwicklung
	\begin{enumerate}
		\item Erstellung der Da Dictionary-Objekte
		\item Erweiterung der /SAPAPO/ORDFLDS
		\item RFC fähiges Funktionsmodul erstellen mit Speicherlogik zur Übertragung der Kommentare   aus dem ECC ins APO
	\end{enumerate}
	\item ECE Entwicklung
	\begin{enumerate}
		\item In den vorhandene Speicherlogiken den RFC einbauen
	\end{enumerate}
	\item SCE Entwicklung
	\begin{enumerate}
		\item Erstellung der Klassen ALV, GUI, Controller und Constants
		\item GUI Funktionsgruppe erstellen
		\item Maintenance Screen erstellen
		\item RRP3 Enhancement implementieren
	\end{enumerate}
\end{enumerate}
\clearpage

\subsection{Use Case-Diagramm}
\label{app:UseCase}
\begin{figure}[htb]
\centering
\includegraphicsKeepAspectRatio{UseCase.pdf}{0.7}
\caption{Use Case-Diagramm}
\end{figure}

\subsection{Pflichtenheft (Auszug)}
\label{app:Pflichtenheft}

\subsubsection*{Zielbestimmung}

\begin{enumerate}[itemsep=0em,partopsep=0em,parsep=0em,topsep=0em]
\item Musskriterien % Wikipedia: für das Produkt unabdingbare Leistungen, die in jedem Fall erfüllt werden müssen
	\begin{enumerate}
	\item ECC System
		\begin{itemize}
		\item Ein neues Package muss im ECC angelegt werden mit dem Namen /CAMELOT/OC.
		\item Die AUFK Datenbanktabelle muss um ein Feld mit dem Namen ZZ/\_ORDER/\_COM-MENT erweitert werden
		\item Ein extra Programm welches eine gefilterte Auswahl an Prozess-, Plan- und Produktions-Aufträgen anzeigt muss erstellt werden.
		\item Ebenfalls sollen hier jeweils die Kommentare für die Aufträge angezeigt und vom User geändert werden können. 
		\item Weiterhin soll dieses Programm massenänderungs fähig sein. 
		\item Außerdem soll es ein zweites Programm geben, in welche der User einen Tabellenfeld der AUFK Datenbanktabelle angeben kann, welches dann für die Kommentare genutzt wird.
		\item Mittels Input Checks soll verhindert werden, dass der User ein Feld angibt, dass es in der AUFK Tabelle nicht gibt
		\item Zu aller Letzt soll die COR1-3 erweitert werden.
		\begin{itemize}
			\item Ein neues Package muss angelegt werden mit dem Namen ZPP, in welches dann die COR Erweiterung kommt.
			\item Mittels der cmod Transaktion muss ein neues Projekt mit dem Namen Z\_COR angelegt werden und zwei Erweiterungspunkte hinzugefügt werden (PPCO0007 und PPCO0020).
		\end{itemize}
		\end{itemize} 
	\item APO System. 
		\begin{itemize}
		\item Es muss ebenfalls ein neues Package angelegt werden /CAMELOT/OC
		\item Die /SAPAPO/ORDFLDS muss um ein Feld mit dem Namen Order\_Comment erweitert werden.
		\item Ebenso wie im ECC System soll es ein Programm geben, welches mittels vom User eingegebener Order Number einen Auftrag mit Kommentar anzeigt, welcher vom User geändert werden kann.
		\item Desweiteren soll die RRP3 erweitert werden.
		\begin{itemize}
			\item Es muss ein neues Package mit dem Namen ZSCM angelegt werden.
			\item Mittels der se18 muss ein BAdI in diese Package implementiert werden.
			\item Die Logik muss programmiert werden.
			\item Eine implizite Erweiterung muss angelegt werden, um die T\_style Tabelle zu manipulieren.
		\end{itemize}
		\end{itemize}
	\end{enumerate}
\end{enumerate}


\clearpage
% !TEX root = Projektdokumentation.tex
\subsection{Dictionary-Objekte}
\label{Dictionary-Objekte}

\subsubsection{Tabellentypen}
\label{Tabellentypen}

\tabelleAnhang{Tabellentypen}

\subsubsection{Strukturen}
\label{Strukturen}

\tabelleAnhang{Strukturen}

\subsubsection{Datenelemente}
\label{Datenelemente}

\tabelleAnhang{Datenelemente}

\subsubsection{Domänen}
\label{Domänen}

\tabelleAnhang{Domänen}
\clearpage

\subsection{Screenshots der Anwendung}
\label{Screenshots}
\begin{figure}[htb]
\centering
\includegraphicsKeepAspectRatio{AdministrationScreenECC.png}{1}
\caption{Ändern des Feldnamens für den Kommentar in der AUFK Datenbanktabelle}
\end{figure}
\clearpage
\begin{figure}[htb]
\centering
\includegraphicsKeepAspectRatio{MaintenanceSelScreenECC.png}{1}
\caption{Selektionsbildschirm für den Maintenance Screen}
\end{figure}
\clearpage
\begin{figure}[htb]
\centering
\includegraphicsKeepAspectRatio{MaintenanceALVScreenECC.png}{1}
\caption{}
\end{figure}
\clearpage
\begin{figure}[htb]
\centering
\includegraphicsKeepAspectRatio{MaintenanceSelPopupECC.png}{1}
\caption{}
\end{figure}
\clearpage
\begin{figure}[htb]
\centering
\includegraphicsKeepAspectRatio{MaintenanceSelScreenAPO.png}{1}
\caption{}
\end{figure}
\clearpage
\begin{figure}[htb]
\centering
\includegraphicsKeepAspectRatio{MaintenanceALVScreenAPO.png}{1}
\caption{}
\end{figure}
\clearpage
\begin{figure}[htb]
\centering
\includegraphicsKeepAspectRatio{RRP3DisplayMode.png}{1}
\caption{}
\end{figure}
\clearpage
\begin{figure}[htb]
\centering
\includegraphicsKeepAspectRatio{RRP3EditMode.png}{1}
\caption{}
\end{figure}
\clearpage
\begin{figure}[htb]
\centering
\includegraphicsKeepAspectRatio{COR3.png}{1}
\caption{}
\end{figure}
\clearpage
\begin{figure}[htb]
\centering
\includegraphicsKeepAspectRatio{COR2.png}{1}
\caption{}
\end{figure}
\clearpage
\clearpage
\subsection{programmierrichtlinien}
\label{app:programmierrichtlinien}
Im folgenden Abschnitt werden kurz die am häufigsten gebrauchten programmierrichtlinien des CAMELOT ITLabs aufgeführt.

\subsubsection{Bennenung}:
\paragraph{Variablen}
Globale Variablen z.B. Variablen in Reports, müssen immer mit einem g(global) anfangen. 
Lokale Variablen z.B. in Methoden von Klassen, müssen immer mit einem l(local) anfangen.

Darauf folgt dann entweder ein:
	t für Tabellen
	s für Strukturen
	v für Variablen
	
Darauf folgt dann ein \_ und ein passender Name welcher die Variable möglichst treffsicher beschreibt.

\paragraph{Klassen}
Klassen Bezeichnungen starten immer mit dem Paketnamen in welchem sich die Klassen befindet + CL(Class) z.B. /CAMELOT/CL\_OC. Der Paketname ist /CAMELOT/OC und in die Mitte wird ein CL gehängt.
	


\clearpage

\subsection{Klassendiagramm}
\label{app:Klassendiagramm}
\begin{figure}[htb]
\centering
\includegraphicsKeepAspectRatio{Klassendiagramm.pdf}{1}
\caption{Klassendiagramm}
\end{figure}
\clearpage


\end{document}
